\subsection{Subgrupo transitivo}
\begin{frame}{Subgrupo transitivo}
  \begin{block}{Definição}
    Dado um conjunto $X$, dizemos que um subgrupo $S$ de $Sim(X)$ é um \textbf{subgrupo transitivo} se dados $x,y\in X$, temos que existe $\sigma\in S$ tal que $\sigma(x)=y$. 
  \end{block}
  \begin{block}{Lema}
    Seja $G$ um grupo e $H\leq G$. Definamos $X=\{xH:x\in G\}$ (o conjunto das classes laterais a esquerda de $H$) e $\varphi$ da seguinte maneira
    \begin{align*}
      \varphi:G&\to Sim(X),\\
              g&\to \varphi_g:X\to X,\\
    \phantom{g}&\phantom{\to \varphi_g:}xH\to gxH.
    \end{align*}
    Então $\Img \varphi$ é um subgrupo transitivo de $Sim(X)$.
  \end{block}
\end{frame}
\subsection{$S_6$ não é completo}
\begin{frame}{}
  \begin{block}{Lema}
    Existe um subgrupo transitivo $K$ de $S_6$ com ordem $120$ que não contém transposições. 
  \end{block}
  \begin{block}{Proof}
    \begin{enumerate}
      \item Pelos teoremas de Sylow temos que existe $H\leq S_5$ tal que $[G:H]=6$.\\
      \item Definimos $X=\{aH:a\in S_5\}$, logo $|Sim(X)|\cong S_6$.
      \item 
        \begin{align*}
          \varphi:S_5&\to Sim(X),\\
                  g&\to \varphi_g:X\to X,\\
        \phantom{g}&\phantom{\to \varphi_g:}xH\to gxH.
        \end{align*}
      \item $K\cong \Img\varphi$.
    \end{enumerate}
  \end{block}
\end{frame}
\begin{frame}{}
  \begin{block}{Proof}
    Raciocinemos pela contradição
    \begin{enumerate}
      \item Suponha $\alpha=(12345)\in K$ e $(ij)\in K$ transposição.
      \item Como $K$ é transitivo temos que existe $\beta\in K$ tal que $\beta(j)=6$.
      \item $\beta(ij)\beta^{-1}=(\beta(i),6)$.
      \item Conjugando $(\beta(i),6)$ pelas potencias de $\alpha$ temos que $(16),(26),(36),(46),(56)\in K$.
      \item $\left\langle (16),(26),(36),(46),(56) \right\rangle=S_6$. Absurdo.
    \end{enumerate}        
  \end{block}
\end{frame}
\begin{frame}{}
  \begin{block}{Teorema}
    Existe um automorfismo de $S_6$ que não é interno.     
  \end{block}
  \begin{block}{Proof}
    \begin{enumerate}
      \item $[S_6:K]=6$, considere $X=\{\alpha_i K:i\in\{1,\cdots,6\}\}$.
      \item
        \begin{align*}
          \sigma:S_6&\to Sim(X),\\
                   g&\to \sigma_g:X\to X,\\
                    &\phantom{\to\sigma_g}\,\alpha_i K\to g\alpha_i K.
        \end{align*}
      \item Logo $\sigma=Aut(S_6)$.
    \end{enumerate}
  \end{block}
\end{frame}
\begin{frame}
  \begin{block}{}
    Raciocinemos pela contradição
    \begin{enumerate}
      \item Suponha que $\sigma\in Inn(S_6)$, logo $\sigma[(12)]=\sigma_{(12)}$ é uma transposição, assim
      \item
        \begin{align*}
          \sigma:S_6&\to S_6,\\
                   (12)&\to \sigma_{(12)}:X\to X,\\
                    &\phantom{\to\sigma_g}\,\alpha_i K\to (12)\alpha_i K.
        \end{align*}
      \item Existe $j$ tal que $(12)\in \alpha_{j}K$, logo $(12)\alpha_jK=\alpha_j K$.
      \item Logo $\alpha_j^{-1}(12)\alpha_j\in K$. Absurdo.
    \end{enumerate}
  \end{block}
  \begin{block}{Teorema}
    $S_6$ não é completo.        
  \end{block}
\end{frame}
